\section{Exceptional divisors between trace models}
\label{sec:trace-exceptional-divisors}

We return to the trace data fixed at the beginning of the chapter, with $\dim Y=r$, and write $S_Z=S_Z(u;Y)$. By \cref{cor:countable-divisorial-pullback}, the full trace current can now be pulled back under a modification $g\colon Z'\to Z$. The pullback is bounded above by the trace current on $Z'$, with equal coefficients along strict transforms, so the difference is a finite effective divisor supported on the exceptional locus. Its logarithmic potential measures the change in the ordinary singularity types of the divisorial parts.

\begin{theorem}\label{thm:trace-exceptional-difference}\index{Exceptional divisor!trace correction and cocycle}
For a smooth compact modification $g\colon Z'\to Z$ over $Y$, there is a unique finite effective real divisor $E_g$ supported on the exceptional locus of $g$ such that
\begin{equation}\label{eq:trace-exceptional-difference}
 S_{Z'}=g^*S_Z+[E_g].
\end{equation}
For a further smooth compact modification $h\colon Z''\to Z'$, these divisors satisfy
\begin{equation}\label{eq:trace-exceptional-cocycle}
 E_{g\circ h}=E_h+h^*E_g.
\end{equation}
Their coefficients are additive and positively homogeneous in the ambient potential $u$.

For every $m$, the analytic-ideal divisors on the two models differ by a finite effective divisor $E_{m,g}$ supported on the exceptional locus:
\[
 [E_{m,g}]=D_{m,Z'}-g^*D_{m,Z}.
\]
\Needspace{5\baselineskip}
For every Gauduchon metric $\gamma'$ on $Z'$, the measures
\[
 [E_{m,g}]\wedge(\gamma')^{r-1}
     \longrightarrow [E_g]\wedge(\gamma')^{r-1}
\]
converge in total variation.
\end{theorem}
\begin{proof}
\emph{Step 1: the exceptional divisor.} Enumerate the prime divisors on $Z$ with positive trace coefficient. By \cref{thm:trace-divisorial-model-compatibility}, the pullback of every finite partial sum of $S_Z$ is at most $S_{Z'}$, and \cref{cor:countable-divisorial-pullback} identifies the increasing limit of these pullbacks with $g^*S_Z$. Thus $S_{Z'}-g^*S_Z$ is a positive closed divisorial current. The two currents have equal coefficients along the strict transform of every prime divisor on $Z$, because these coefficients are computed by the same valuation. Every other prime divisor on $Z'$ is exceptional for $g$, and the exceptional analytic set has only finitely many irreducible components on the compact manifold $Z'$. Hence the difference is the integration current of a finite effective real divisor $E_g$ supported on its divisorial components, and the coefficients of integration currents determine $E_g$ uniquely.

\emph{Step 2: the cocycle identity.} Pullback of positive closed $(1,1)$-currents is functorial, since it is defined by composing local plurisubharmonic potentials. Using \eqref{eq:trace-exceptional-difference} for $h$ and for $g$ gives
\[
 S_{Z''}=h^*S_{Z'}+[E_h]
        =(g\circ h)^*S_Z+[h^*E_g]+[E_h],
\]
and uniqueness gives \eqref{eq:trace-exceptional-cocycle}. Additivity and positive homogeneity follow from the same properties of the trace coefficients and the linearity of pullback.

\Needspace{10\baselineskip}
\emph{Step 3: approximation.} For fixed $m$, apply \eqref{eq:trace-ideal-model-inequality} to the finitely many prime divisors in the support of the image ideal on $Z$. After division by $m$, it gives $D_{m,Z'}-g^*D_{m,Z}\geq0$. The coefficients along strict transforms vanish, so this difference is the integration current of an effective real divisor $E_{m,g}$ supported on the exceptional locus.

Choose also a Gauduchon metric $\gamma$ on $Z$. By \cref{thm:trace-divisorial-current-budget}, the currents $S_Z-D_{m,Z}$ and $S_{Z'}-D_{m,Z'}$ are positive, and their trace masses tend to zero. Subtracting the two identities for the exceptional currents gives
\[
 [E_{m,g}]-[E_g]
   =-(S_{Z'}-D_{m,Z'})+g^*(S_Z-D_{m,Z}),
\]
where both $S_{Z'}-D_{m,Z'}$ and $g^*(S_Z-D_{m,Z})$ are positive. By \eqref{eq:modification-linear-mass-bound}, the total variation of the signed measure $([E_{m,g}]-[E_g])\wedge(\gamma')^{r-1}$ is therefore at most
\[
 \int_{Z'}(S_{Z'}-D_{m,Z'})\wedge(\gamma')^{r-1}
 +C_g\int_Z(S_Z-D_{m,Z})\wedge\gamma^{r-1}.
\]
Both terms tend to zero, which proves convergence in total variation.
\end{proof}

For a modification $g\colon Z'\to Z$ as in the theorem, the coefficient of $E_g$ along an exceptional prime divisor $F\subset Z'$ is
\begin{equation}\label{eq:trace-exceptional-coefficient}
 \tau_{Y,u}(F)
  -\sum_{D\subset Z}\tau_{Y,u}(D)\operatorname{ord}_F(g^*D),
\end{equation}
where $D$ ranges over the prime divisors on $Z$ with positive trace coefficient. This countable sum converges: it is the coefficient of $g^*S_Z$ along $F$ in the positive divisorial series of \cref{cor:countable-divisorial-pullback}, and it is at most $\tau_{Y,u}(F)$. The same current identity compares the ordinary singularity types attached to $S_Z$ and $S_{Z'}$ by \cref{cor:divisorial-current-ordinary-type}.

\Needspace{14\baselineskip}
\begin{corollary}\label{cor:trace-model-ordinary-comparison}\index{Divisorial current!ordinary types under pullback}
Choose decompositions $S_Z=B_Z+\ddc v_Z$ and $S_{Z'}=B_{Z'}+\ddc v_{Z'}$ with smooth closed forms $B_Z,B_{Z'}$, as in \cref{cor:divisorial-current-ordinary-type}. Let $\rho_{E_g}$ be a global logarithmic potential of $E_g$, that is, a quasi-plurisubharmonic function on $Z'$ such that $\ddc\rho_{E_g}-[E_g]$ is smooth. Then
\[
 v_{Z'}-(v_Z\circ g)-\rho_{E_g}
\]
is represented by a smooth function. In particular,
\[
 v_{Z'}\preceq v_Z\circ g
\]
in the sense of \cref{def:singularity-order}, and the two ordinary singularity types are equal if and only if $E_g=0$.
\end{corollary}
\begin{proof}
Global logarithmic potentials exist: choose smooth Hermitian metrics on the line bundles of the prime components of $E_g$, let $\alpha_{E_g}$ be the corresponding weighted sum of their curvature forms, and let $\rho_{E_g}$ be the weighted sum of the logarithms of the squared norms of the canonical sections. Then
\[
 \ddc\rho_{E_g}=[E_g]-\alpha_{E_g},
\]
and any two global logarithmic potentials differ by a smooth function, by \cref{lem:closed-current-global-potential}. We may therefore use this choice. The function $\rho_{E_g}$ is quasi-plurisubharmonic and bounded above on the compact manifold $Z'$.

\Needspace{10\baselineskip}
The chosen decompositions and \eqref{eq:trace-exceptional-difference} give
\[
 \ddc\bigl(v_{Z'}-v_Z\circ g-\rho_{E_g}\bigr)
       =g^*B_Z+\alpha_{E_g}-B_{Z'},
\]
which is smooth. By \cref{lem:closed-current-global-potential}, this difference of locally integrable functions is represented by a smooth function. The resulting equality of quasi-plurisubharmonic functions extends from almost everywhere to everywhere by \cref{lem:qpsh-almost-everywhere-order}. Since $\rho_{E_g}$ is bounded above, this gives the ordinary domination.

If $E_g=0$, the smooth difference gives ordinary equivalence. Conversely, ordinary equivalence would give equal generic Lelong numbers along every prime divisor of $Z'$, whereas the identity just proved shows that along a prime component of $E_g$ these numbers differ by its positive coefficient. Hence ordinary equivalence forces $E_g=0$.
\end{proof}

When the restricted background $\theta|_Y$ is closed, the positive residual currents $R_Z$ constructed in the proof of \cref{prop:trace-bott-chern-bclass}, and their Bott--Chern classes, satisfy
\[
 \begin{split}
 R_{Z'}&=g^*R_Z-[E_g],\\
 \mathbb A_{Z'}(u;Y,\theta)
       &=g^*\mathbb A_Z(u;Y,\theta)-\{[E_g]\}_{\mathrm{BC}}.
 \end{split}
\]
Indeed, the smooth background and the common trace potential pull back, while \eqref{eq:trace-exceptional-difference} gives the change in the subtracted divisorial current. The correction $\{[E_g]\}_{\mathrm{BC}}$ accounts for the additional divisorial vanishing prescribed on the higher model. This is consistent with the exact pullback of the trace envelope in \cref{thm:trace-envelope-modification}, since the envelope imposes the trace inequalities on all models at once. The following example shows that the correction can be nonzero even when the trace data are realized exactly by a potential.

\Needspace{9\baselineskip}
\begin{example}\label{ex:exceptional-trace-with-exact-potential}
Let $Y=\mathbb P^2$ with homogeneous coordinates $[t_0:t_1:t_2]$, let $\sigma$ be the Fubini--Study form normalized by $\int_Y\sigma^2=1$, and choose $a\in\mathbb Q_{>0}$. Put
\[
 u=a\log\frac{|t_1|^2+|t_2|^2}
                   {|t_0|^2+|t_1|^2+|t_2|^2},
\]
and take the ambient manifold to be $Y$ itself. Then $u\in\PSH(Y,a\sigma)$. It has analytic singularities, with its only pole at $p=[1:0:0]$, and zero generic Lelong number on $Y$. By \cref{prop:trace-naive-comparison}, its trace data are given by the generic Lelong numbers of its pullbacks: $\tau_{Y,u}(E)=\nu(u\circ\pi,E)$ for every prime divisor $E$ on a smooth compact modification $\pi$ of $Y$. Since $u$ has no divisorial poles on $Y$, we have $S_Y=0$, whose canonical ordinary type is that of a smooth function. If $\pi\colon\widetilde Y\to Y$ is the blowup at $p$ with exceptional curve $F$, then
\[
 S_{\widetilde Y}=a[F],\qquad E_\pi=aF.
\]
Indeed, on either standard chart of the blowup, the numerator is the squared modulus of the exceptional coordinate times a smooth positive function, so the pole order is $a$ along $F$ and zero along every other prime divisor. Thus the trace data are realized exactly by $u$, although $E_\pi\neq0$. The exceptional difference does not obstruct exact realization; it prevents the identification of the realizing potential with the canonical ordinary type of the divisorial part on the original model, which here is smooth.

\Needspace{8\baselineskip}
The same phenomenon occurs on a manifold without the bounded mass property. Let $(W,\omega_W)$ be a connected compact Hermitian manifold of dimension $k$ that fails the bounded mass property, such as the Hopf threefold of \cref{thm:hopf-unbounded}, and pull $u$ back to $Y\times W$; take both the ambient manifold and the submanifold to be $Y\times W$. The pole set of the pullback is $\{p\}\times W$, of codimension two, and the modification $\widetilde Y\times W\to Y\times W$ has exceptional coefficient $a$ along $F\times W$. Since the pullback of $u$ has analytic singularities, \cref{prop:trace-naive-comparison} again shows that it realizes the trace data exactly. Let $\operatorname{pr}_Y$ and $\operatorname{pr}_W$ be the two product projections.

\Needspace{9\baselineskip}
Choose smooth $\omega_W$-plurisubharmonic functions $H_j$ whose masses tend to infinity. For the Hermitian form $\Omega=\operatorname{pr}_Y^*\sigma+\operatorname{pr}_W^*\omega_W$, the pullbacks of the $H_j$ have masses
\[
 \binom{k+2}{2}\left(\int_Y\sigma^2\right)
          \left(\int_W(\omega_W+\ddc H_j)^k\right)\longrightarrow+\infty,
\]
by the binomial expansion of smooth forms and the dimensions of the two factors. Hence the product also fails the bounded mass property. The pulled-back potential belongs to $\PSH(Y\times W,C\Omega)$ for every $C\geq\max\{a,1\}$, since its curvature is bounded below by $-a\operatorname{pr}_Y^*\sigma$. Scaling the smooth functions $H_j$ by $C$ shows that $C\Omega$ also has unbounded mass.
\end{example}
